\documentclass{article}

% \usepackage{neurips_2024} % ready for submission
\usepackage[preprint]{neurips_2024} % to compile a preprint version, e.g., for submission to arXiv, add the [preprint] option
% \usepackage[final]{neurips_2024} % to compile a camera-ready version, add the [final] option
% \usepackage[nonatbib]{neurips_2024} % to avoid loading the natbib package, add option nonatbib

% === INÍCIO DAS ANOTAÇÕES ADVINDAS DO MODELO NEURIPS_2024 === %

% \PassOptionsToPackage{options}{natbib} % If you wish to load the \verb+natbib+ package with options, you may add the following before loading the \verb+neurips_2024+ package:
% \citet{fulano_2018} afirmam que \dots
% \footnote{Sample of the first footnote.}
% Caption after the figure
% Caption before table
% \begin{ack} % Agradecimento aos institutos financiadores e/ou empresas envolvidas no projeto. <https://neurips.cc/Conferences/2024/PaperInformation/FundingDisclosure>
% \end{ack}
% referências podem ser \small (9 pts); qualquer estilo de citação, desde que consistente.
% \appendix % Optionally include supplemental material (complete proofs, additional experiments and plots) in appendix. All such materials \textbf{SHOULD be included in the main submission.}

% === FIM DAS ANOTAÇÕES ADVINDAS DO MODELO NEURIPS_2024 === %

\usepackage[brazil]{babel}      % Syllabification and hyphenation for Brazilian Portuguese
\usepackage[inline]{enumitem}   % inline enumeration
\usepackage[T1]{fontenc}        % use 8-bit T1 fonts
\usepackage[utf8]{inputenc}     % allow utf-8 input
\usepackage{amsfonts}           % blackboard math symbols
\usepackage{booktabs}           % professional-quality tables
\usepackage{comment}            % enables the comment environment
\usepackage{hyperref}           % hyperlinks
\usepackage{microtype}          % microtypography
\usepackage{multicol}           % multiple columns
\usepackage{nicefrac}           % compact symbols for 1/2, etc.
\usepackage{url}                % simple URL typesetting
\usepackage{xcolor}             % colors

\title{Project Report\\JSONLLM: Extração Escalável de Atributos de Produtos com Grandes Modelos de Linguagem (LLMs)}

\author{
    João Vítor Fernandes Dias
    % \thanks{ORCID: \href{https://orcid.org/0000-0002-8156-9551}{0000-0002-8156-9551}}
    \thanks{
        ORCID: \href{https://orcid.org/0000-0002-8156-9551}{0000-0002-8156-9551};
        GitHub: \href{https://github.com/jvfd3}{jvfd3};
        LinkedIn: \href{https://www.linkedin.com/in/jvfd3}{jvfd3}
    } \\
    Computer Science Postgraduate Program\\
    2024711370\\
    Federal University of Minas Gerais\\
    Belo Horizonte, MG; 31270-901 \\
    \texttt{joaovitorfd2000@gmail.com} \\
    \And % Quebra de linha automática
    % \AND % Quebra de linha forçada
    Melissa Dias Mattos \\
    Computer Science Postgraduate Program\\
    2024675063\\
    Federal University of Minas Gerais\\
    Belo Horizonte, MG; 31270-901 \\
    \texttt{melissadmattos366@gmail.com} \\
}

\begin{comment}
\appendix
\section{Appendix / supplemental material}
\newpage
\section*{NeurIPS Paper Checklist}
\end{comment}
